\documentclass[12pt]{article}
\usepackage[utf8]{inputenc}
\usepackage[T1]{fontenc}
\usepackage{hyperref}
\title{Wakabunga Language Resources for K-12 Curriculum Development}
\author{Compiled by Claude Sonnet 5.5 (claude-sonnet-5-5) for Andr\'as Kornai}
\date{}
\begin{document}
\maketitle

\section{Introduction}
Wakabunga (ISO 639-3: wwb) was an Aboriginal language of Queensland, Australia. It has been extinct since about the 1960s, and its classification is uncertain. Wikipedia reports that the main word list labelled ``Wakabunga'' is regarded by other sources as Kalkatungu, so it may not represent Wakabunga at all. Barry Blake suggested that a different word list, compiled by Montagu Curr, reflects a Mayi dialect, which would make Wakabunga a possible Mayi language. The language is probably unattested, so no reliable material exists for teaching it.

\section{Primary Educational Resources}
None found. With no reliable linguistic record, a language-teaching curriculum is not currently possible.

\section{Linguistic and Reference Materials}
\begin{itemize}
\item Wikipedia entry on the language, with the classification discussion: \url{https://en.wikipedia.org/wiki/Wakabunga_language}.
\item Wikipedia entry on the Wakabunga people: \url{https://en.wikipedia.org/wiki/Wakabunga}.
\item South Australian Museum archives page on the Wakabunga (Queensland) language group, based on Norman Tindale's records: \url{https://www.samuseum.sa.gov.au/collection/archives/language_groups/wakabunga}. Tindale placed their lands at about 13,000 km$^2$ around the upper Leichhardt River and Gunpowder Creek and linked them to the Kalkatungu.
\end{itemize}

\section{Religious and Community Resources}
No scripture or community-organisation resources were found.

\section{Digital and Audio Resources}
None found.

\section{Curriculum Development Context}
A heritage curriculum cannot draw on Wakabunga language material. A realistic option is a place-based unit on the people, country and history of the upper Leichhardt River area. It would be developed with the Traditional Owners' descendants and with the neighbouring Kalkatungu and Mayi communities. Any use of the Kalkatungu or Mayi languages would need permission from those communities. The unresolved status of the word lists should be stated plainly in any teaching material.

\end{document}